大语言模型的能力由参数规模 $N$、数据量 $D$、数据质量 $Q$ 与领域配比 $\bp$ 共同决定，而扩大规模、提升质量与延长上下文都要消耗有限的算力。本文沿赛题“数据刻画—规律建模—资源优化—演进预测”的递进主线建立四个模型，以接口文件组织各问结果，并按证据等级（来源观测、半合成、外推估算）区分各附件的用途。

针对问题一，将 22 个质量字段展开为 25 个方向统一的指标，以 CRITIC 权重和加权 Huber 中心构造样本级主质量分 $Q$，并取中位数聚合为领域质量分。A1 与对应全量扩展集 A2/A3 的分布未检出显著差异（KS 检验 $p=0.313/0.220$）。评分后以六个来源定义高分歧候选；A1/A2/A3 的候选比例为 $5.00\%/4.73\%/4.99\%$，多来源对立候选比例为 $1.36\%/1.58\%/1.77\%$。基于 A16/A18 将七个质量领域桥接到 17 个配比领域，并用 clr 回归建立配比与 13 个验证域交叉熵的定量关系；1M 留出集的平均逐域 $R^2=0.772$。GBDT 搜索得到的等权候选配比相对均匀配比预测降损 $9.915\%$。冻结 1M 模型在 60M/1B 留出集的配比排序 Spearman 为 $0.467/0.677$，但绝对损失预测不具跨尺度可迁移性；A12--A15 为估算/外推数据，仅用于边界讨论。

针对问题二，先以 B1 校准经典标度律，再用 B2、B4、B5 检查跨来源趋势。固定经典基线后，在 B6 拟合广义质量项 $L=L_0+(1-Q_B)(\kappa_NAN^{-\alpha}+\kappa_DBD^{-\beta}+k_0)$，并在 B7 新增 90 个半合成点上留出检验（RMSE $0.042$）；B8 的质量方向与 B6/B7 不一致，不并入拟合。A、B 之间的映射在聚合可加及 B 侧惩罚仿射等假设下可推出为仿射，但锚点仍是建模选择；配比组成修正尚未标定。该半合成模型给出质量弹性、质量地板、质量—规模替代条件及算力最优规模的解析式，结论不等同于真实训练规律。

针对问题三，三项成本都与 $D$ 成正比，最优时预算约束必取等，据此把 $D$ 解析消去，将问题化为 $(\log N,Q)$ 上的二维无约束问题，并用“内层黄金分割 + 外层三级网格加密”的向量化求解器，在 3 种成本函数 × 7 档预算 × 5 档 $L_{\mathrm{ctx}}$ × 2 种基线质量下求得 210 组最优解。数值解与闭式解的相对误差不超过 $1.83\times10^{-7}$，与穷举网格的 Loss 差不超过 $1.6\times10^{-5}$。以 KKT 活跃约束集的改变定义结构性转移，并由边际比 $\Phi(C)=1$ 解析识别临界预算。旧接口主口径（$Q_0=0.988$，$L_{\mathrm{ctx}}=2048$）下，$\log_{10}C_{\mathrm{crit}}$ 为：指数型 $20.61$、幂函数型 $20.48$、对数型 $18.52$。三种识别口径的极差不超过 $0.05\dex$，bootstrap 90\% 区间在 $\pm0.03\dex$ 以内。转移几乎呈“不提质 → 提满”的跳变，210 组中 34 组为 $Q_0$ 角点、3 组为内点，均出现在 $10^{19}$–$10^{20}$。解析给出 $L_{\mathrm{ctx}}^{\mathrm{crit}}=6/\eta=30000$；上下文增至 131k 时，$C=10^{22}$ 的最优模型缩至 $0.48$ 倍，注意力成本占 $81\%$。将配比联立优化可再降 $13.7\%$，但几乎全部来自配比效应；最优配比排序在不同预算下的秩相关 $\ge0.976$，“先定配比、再定规模”的分离策略成立。

针对问题四，能力的主口径为 C1 的 Average，并对 C8 的 39 个叶子任务逐一聚合，识别出 10 个饱和任务，构造去饱和综合分，发现饱和任务使前沿增速被低估约 1.6 倍。将 base 与 chat 分开处理，时间轴采用 C2 发布日。按可比性等级分层，用加权 sigmoid 拟合 C6 的 Loss–Benchmark 桥接（$r=-0.51$，留一 RMSE $8.08$，优于线性的 $8.28$）。以问题二的标度律 $\hat L(N,D)$ 作为规模状态，建立“S 形桥接 + 时间技术项”的结构模型（$R^2=0.87$），并用 Shapley 分解前沿增长。2023.5→2025.0 期间 base 前沿由 $12.5$ 升至 $37.3$，结构口径下非规模技术进步占 $13\%$–$23\%$，线性效率前沿口径为 $38\%$–$46\%$；两类方法的差异来自对 S 形加速段的归因不同。结论是规模扩张占主导。预测以问题三的最优配置为算力通道，在基线 / 减半 / 1/4 三种算力增速（$0.60/0.30/0.15\dex/\mathrm{yr}$）情景下，base 开源前沿 24 个月中位数为 $46.2/45.2/44.1$。算力放缓时，技术进步在增量中的占比由 $55\%$ 升至 $77\%$。回测误差为 $-0.78$，而朴素线性外推为 $+2.65$。

问题二的质量形式与问题一当前质量接口尚未贯通，问题三、四的上述数值仍基于旧的 $Q$ 锚点和旧标度律接口；在新接口下完成重估前，这些下游数值不应解读为修订模型的预测结果。

各问按数据结构分别采用留出检验、跨来源趋势对照或解析—数值核查，《数据说明》中的“建议”也逐条核对。B6/B7 的新质量项目前报告点估计与半合成留出结果，未完成可复现的区间估计；问题三、四仍使用旧接口，需贯通重算后再解释为修订模型下的结论。半合成与外推数据只支持其明确标注的形式与边界讨论。

\vspace{6pt}
\noindent{\heitibf 关键词：}CRITIC 赋权；成分数据 clr 回归；广义标度律；KKT 结构性转移；Shapley 贡献分解
